% Insert after the existing appendices; do not issue another \appendix.
% Requires amsmath and amssymb. Add bi_hmc_references.bib to the bibliography.
% Settings checked against the saved Python Monte Carlo/record-length results.
% The conditional preamble supports the editor preview of this same file;
% it is skipped when the appendix is included in an existing manuscript.
\ifdefined\section
  \let\BIHMCfinishstandalone\relax
\else
  \documentclass[11pt]{article}
  \usepackage[a4paper,margin=25mm]{geometry}
  \usepackage{amsmath,amssymb}
  \def\BIHMCfinishstandalone{%
    \bibliographystyle{plain}%
    \bibliography{bi_hmc_references}%
    \end{document}}
  \begin{document}
  \appendix
\fi

\section{Bayesian identification using Hamiltonian Monte Carlo}
\label{app:bi_hmc}

The Bayesian baseline, denoted BI-HMC, estimates the same physical
coefficient vector $\theta\in\mathbb R^n$ as the other methods, together
with the innovation variance $\sigma_\nu^2$. The model orders and the
normalization $f_1=1$ are retained. Assuming
$\nu(t)\overset{\mathrm{iid}}{\sim}\mathcal N(0,\sigma_\nu^2)$,
independently of the input, Bayes' rule gives
\begin{equation}
\label{eq:bi_posterior}
 p(\theta,\sigma_\nu^2\mid C,u)
 \propto p(C\mid\theta,\sigma_\nu^2,u)\,
          p(\theta)\,p(\sigma_\nu^2).
\end{equation}
The deterministic signal $\beta(t;\theta)$ is evaluated through the
original IIR recursion and polynomial nonlinearity. Thus BI-HMC provides
one common reference for the PL and FSM comparisons, without an FSM
truncation in its likelihood.

\paragraph{Likelihood.}
Conditionally on the measured input $u$ and a candidate $\theta$, the
signal $\beta(t;\theta)$ is deterministic. The only remaining randomness
in $c(t)$ is the additive AR disturbance, which is a linear filtering of
Gaussian innovations. In particular, let $B$ denote the number of burn-in
samples and collect all innovations affecting the retained record in
$V_B=[\nu(1-B),\ldots,\nu(\tau)]^\top$. With a zero noise state before
these samples,
\begin{equation}
\label{eq:bi_gaussian_construction}
 \begin{aligned}
 V_B&\sim\mathcal N(0,\sigma_\nu^2 I_{\tau+B}),\\
 C&=\boldsymbol\beta(\theta)+H_B(d)V_B,
 \end{aligned}
\end{equation}
where $\boldsymbol\beta(\theta)=[\beta(1;\theta),\ldots,
\beta(\tau;\theta)]^\top$ and $H_B(d)$ is the
$\tau\times(\tau+B)$ convolution matrix of the candidate AR noise
filter, restricted to the retained output samples. An affine
transformation of a Gaussian vector is Gaussian; hence
\begin{equation}
\label{eq:bi_gaussian_data}
 C\mid\theta,\sigma_\nu^2,u
 \sim\mathcal N\!\left(\boldsymbol\beta(\theta),
                      \sigma_\nu^2 R_d\right),
 \qquad R_d=H_B(d)H_B(d)^\top.
\end{equation}
Thus the mean depends on the signal-path coefficients, while the
covariance depends on the AR coefficients and innovation variance.
The output samples are generally correlated. The nonlinear dependence
of $\boldsymbol\beta$ on $\theta$ does not invalidate this
\emph{conditional} Gaussian distribution or imply a Gaussian posterior.

Define $e(t;\theta)=c(t)-\beta(t;\theta)$ and
$E(\theta)=C-\boldsymbol\beta(\theta)$. Substitution into the
multivariate Gaussian density gives the likelihood
\begin{align}
\label{eq:bi_likelihood}
 p(C\mid\theta,\sigma_\nu^2,u)
 &=\frac{\exp[-Q(\theta)/(2\sigma_\nu^2)]}
 {(2\pi\sigma_\nu^2)^{\tau/2}|R_d|^{1/2}},\\
\label{eq:bi_mahalanobis}
 Q(\theta)&=E(\theta)^\top R_d^{-1}E(\theta).
\end{align}
Taking logarithms yields
\begin{equation}
\label{eq:bi_log_likelihood}
 \begin{aligned}
 \log p(C\mid\theta,\sigma_\nu^2,u)
 &=-\frac{\tau}{2}\log(2\pi)
   -\frac{\tau}{2}\log\sigma_\nu^2\\
 &\quad-\frac12\log|R_d|
   -\frac{Q(\theta)}{2\sigma_\nu^2}.
 \end{aligned}
\end{equation}
Only the first term is constant with respect to all inference
parameters; the covariance-determinant term must also be retained.

The implementation avoids forming or inverting the full $R_d$.
The AR recursion reconstructs
\begin{equation}
\label{eq:bi_residuals}
 \hat\nu(t;\theta)=e(t;\theta)
       +\sum_{i=1}^{n_d}d_i e(t-i;\theta),\qquad t>n_d.
\end{equation}
Let $m=\min(n_d,\tau)$, let $E_m(\theta)$ contain the first $m$
residuals, and let $\Gamma_d$ be their model covariance divided by
$\sigma_\nu^2$. The joint density factors into the initial residual
density and the subsequent independent Gaussian innovation densities.
Since the residual-to-innovation transformation is triangular with
unit diagonal, this gives
\begin{equation}
\label{eq:bi_Q}
 \begin{aligned}
 Q(\theta)
 &=E_m(\theta)^\top\Gamma_d^{-1}E_m(\theta)
   +\sum_{t=m+1}^{\tau}\hat\nu(t;\theta)^2,\\
 |R_d|&=|\Gamma_d|.
 \end{aligned}
\end{equation}
In the simulations, the filters start from zero before the $100$
discarded samples. The observed input prehistory is replayed for each
candidate signal model, and $\Gamma_d$ is computed from the candidate
AR filter's impulse response and this burn-in length. Unobserved noise
prehistory is integrated out; neither true innovations nor true internal
signals enter inference. This initial-density contribution distinguishes
the exact likelihood from a conditional criterion that simply drops the
first $n_d$ samples. For $n_d=0$, the initial quadratic term is zero,
$R_d=I_\tau$, and $|\Gamma_d|=1$.

\paragraph{Priors.}
Zero-mean Gaussian prior kernels are assigned to the physical coefficients:
$a_i,b_i,f_j\sim\mathcal N(0,1)$ for $j\geq2$, and
$d_i\sim\mathcal N(0,0.5^2)$. Their product is restricted and normalized
over the region where both recursive filters are strictly stable.
Independently,
\begin{equation}
\label{eq:bi_variance_prior}
 \sigma_\nu^2\sim\mathrm{IG}(3,0.002),\qquad
 p(\sigma_\nu^2)\propto
 (\sigma_\nu^2)^{-4}\exp\!\left(-\frac{0.002}{\sigma_\nu^2}\right).
\end{equation}
The prior mean is $0.001$, whereas the nominal benchmark innovation
variance is $0.01$. These fixed priors are used across both benchmarks and record
lengths; no sparsity hierarchy is introduced.

\paragraph{Smooth sampling coordinates.}
Stability is enforced through reflection coefficients
\cite{Monahan1984}. For either denominator order $p\in\{n_a,n_d\}$,
introduce unrestricted coordinates $r_j$ and set $\kappa_j=\tanh r_j$.
Starting with an empty coefficient vector, the step-up recursion is
\begin{equation}
\label{eq:bi_schur}
 v_j^{(j)}=\kappa_j,\qquad
 v_i^{(j)}=v_i^{(j-1)}+\kappa_j v_{j-i}^{(j-1)},
 \quad i=1,\ldots,j-1,
\end{equation}
for $j=1,\ldots,p$. The final coefficients $v_i^{(p)}$ give $a_i$
or $d_i$, respectively. The coefficients $b_i,f_j$ remain unchanged,
and $\sigma_\nu^2=\exp(s)$ with $s\in\mathbb R$.
Writing this transformation as
$T(\vartheta)=(\theta,\sigma_\nu^2)$, the sampling log density is
\begin{equation}
\label{eq:bi_sampling_density}
 \log\pi(\vartheta)
 =\log p\!\left(T(\vartheta)\mid C,u\right)
   +\log|\det DT(\vartheta)|.
\end{equation}
The Jacobian includes both reflection-coefficient maps and the variance
factor $\exp(s)$. It preserves the stated physical priors under the
change of coordinates.

\paragraph{Sampling and reported estimate.}
Hamiltonian Monte Carlo (HMC) \cite{Neal2011} constructs Markov chains
whose target density is $\pi(\vartheta)$ in
\eqref{eq:bi_sampling_density}. At each transition, an auxiliary Gaussian
momentum is drawn and the negative log density is treated as a potential
energy. A leapfrog integrator uses its gradient to propose a joint move
of all $n+1$ unknowns through the parameter space. A Metropolis
acceptance step corrects the numerical integration error; on rejection,
the chain remains at its previous state. Repeating these transitions
produces correlated draws representing the posterior distribution,
subject to adequate exploration of its relevant regions.

Four chains are initialized at perturbed versions of the transformed
study starting point, with initial variance $0.01$ and coordinate
perturbation standard deviation $0.1$. A preliminary posterior
maximization, limited to $30$ L-BFGS iterations, provides starting states
for sampling. Each chain then undergoes $500$ warm-up transitions.
This phase allows movement away from the initialization and tunes the
sampler to the posterior scales: the step size is adapted by dual
averaging towards acceptance $0.8$ \cite{HoffmanGelman2014}, while a
diagonal mass matrix adjusts the relative movement of the coordinates.
The initial step-size search starts at $0.01$, and each trajectory uses
a uniformly selected number of leapfrog steps from $\{5,\ldots,10\}$.
Warm-up draws are discarded. After warm-up, the step size and mass
matrix are fixed and $1000$ further draws are retained from each chain,
including repeated states after rejection.

The resulting $4000$ draws are assessed using agreement between chains
and effective sample sizes, which account for their serial dependence.
The configured thresholds are rank-normalized/folded split
$\widehat R\leq1.01$ and bulk and tail effective sample sizes of at
least $400$ for every coefficient and the innovation variance
\cite{VehtariEtAl2021}. Acceptance probabilities and divergent
transitions are also monitored. Diagnostic warnings are reported
separately; the default finite-output summaries retain finite, stable
point estimates even when these diagnostics raise a warning.

Finally, each retained draw is transformed to the physical coefficients
and variance. Denoting these draws from chain $j$ by
$\theta^{(j,\ell)}$ and $(\sigma_\nu^2)^{(j,\ell)}$, the reported
estimates approximate the posterior expectations:
\begin{equation}
\label{eq:bi_posterior_means}
 \begin{aligned}
 \hat\theta_{\mathrm{BI\text{-}HMC}}
 &=\frac{1}{4000}\sum_{j=1}^{4}\sum_{\ell=1}^{1000}
      \theta^{(j,\ell)},\\
 \widehat{\sigma_\nu^2}
 &=\frac{1}{4000}\sum_{j=1}^{4}\sum_{\ell=1}^{1000}
      (\sigma_\nu^2)^{(j,\ell)}.
 \end{aligned}
\end{equation}
Thus each identification record yields one final coefficient estimate
for computing parameter errors and output reconstructions. Averaging
is performed in physical coordinates, after the nonlinear transformation.
Total wall-clock time includes initialization, warm-up, retained
sampling, diagnostics, and assessment; sampling transitions are not
equated with optimizer iterations. The same saved Bayesian reference
is included in both parametrization reports rather than being fitted
again for PL.
\BIHMCfinishstandalone
